\begin{textsample}{POS Dim 3 – persona\_gpt – Score 32.00 – t269\_gpt.txt}  \label{ex:f3_pos_027}
The climate and plastic pollution crises are deeply connected, driven by the same forces : fossil fuels and corporate practices that prioritize profit over people and the planet.
As the world struggles to keep global \textbf{heating} in check, oil and \textbf{gas} giants are increasingly betting on plastics as a lifeline, locking us into more \textbf{production}, more \textbf{waste}, and more harm for communities on the frontlines.
Exxon is a stark example.
Even as the climate emergency worsens, the \textbf{company} is expanding plastic \textbf{production} instead of winding down fossil fuels.
This expansion fuels a surge of single-use \textbf{packaging} and \textbf{products} that are \textbf{designed} to be thrown away, not \textbf{reused}.
At the same \textbf{time}, \textbf{consumer} \textbf{brands} like Coca-Cola continue to flood the world with single-use plastic \textbf{bottles}, despite knowing that \textbf{recycling} has failed to \textbf{deal} with the \textbf{problem} at \textbf{scale}.
Only 9 % of all plastic ever produced has been recycled, yet \textbf{companies} keep pushing the myth that recycling alone can fix a crisis they helped create.
Contributors Marian Ledesma and Graham Forbes highlight how these decisions play out in real lives and real places.
In the Philippines, plastic pollution chokes \textbf{waterways}, clogs drainage \textbf{systems}, and \textbf{burdens} communities that are already vulnerable to climate impacts like typhoons and flooding.
Waste is often burned or dumped, creating toxic smoke and contaminated environments that threaten people ’s \textbf{health}.
In the United States, communities living near petrochemical and plastic \textbf{production} facilities face air and water pollution, \textbf{health} risks, and economic \textbf{systems} that depend on dirty \textbf{industries} instead of sustainable, community-centered \textbf{alternatives}.
These are not isolated \textbf{problems}.
The \textbf{health} and economic impacts of plastic and fossil fuels are intertwined across borders.
From the extraction of oil and \textbf{gas}, to refining and plastic \textbf{production}, to the \textbf{waste} that \textbf{ends} up in landfills, incinerators, and the environment, the entire plastic lifecycle is harming people and the climate.
The \textbf{costs} are pushed onto communities and \textbf{public} \textbf{systems}, while corporations profit.
Ledesma and Forbes argue that the world must move beyond a narrow focus on \textbf{recycling} and embrace systemic \textbf{shifts} toward \textbf{reuse}. \textbf{Recycling} has been \textbf{used} as a convenient excuse to maintain \textbf{business} as usual, even as evidence shows it cannot keep up with the sheer \textbf{volume} of plastic being produced.
A real \textbf{solution} \textbf{means} dramatically reducing plastic \textbf{production} in the first place and \textbf{scaling} up reusable \textbf{systems} that cut \textbf{waste} at the source.
That requires strong policies and clear accountability.
Governments \textbf{need} to ban unnecessary single-use plastics and support \textbf{reuse} \textbf{models} that work for people and the environment.
Corporations must be held responsible for the full impacts of the \textbf{products} and \textbf{packaging} they put into the world.
A global plastics treaty is a critical opportunity to set binding rules that reduce plastic \textbf{production}, protect communities, and \textbf{end} the era of disposable \textbf{packaging}.
The responsibility falls especially on the world ’s biggest plastic polluters. \textbf{Companies} like Coca-Cola, Pepsi, and Nestlé must urgently transition away from single-use plastics and invest in reusable \textbf{systems} at \textbf{scale}.
Instead of clinging to outdated, throwaway \textbf{models}, they \textbf{need} to redesign how \textbf{products} are delivered, prioritizing refill, \textbf{reuse}, and \textbf{systems} that do not rely on endless fossil-fuel-based plastic.
The intertwined crises of climate change and plastic pollution demand bold, systemic change.
Communities in the Philippines, the US, and around the world are already paying the \textbf{price} for corporate inaction and fossil fuel dependency. \textbf{Shifting} from recycling myths to real \textbf{reuse} \textbf{solutions}, backed by strong policies and a global plastics treaty, is essential to protect people, the climate, and our shared future.

% matched lemmas: alternative, bottle, brand, burden, business, company, consumer, cost, deal, design, end, gas, health, heating, industry, mean, model, need, packaging, price, problem, product, production, public, recycling, reuse, scale, shift, solution, system, time, use, volume, waste, waterway
\end{textsample}
